\begin{frame}{Vurdering af vildt (§49.6)}
    \begin{itemize}
        \item \textbf{Hovedregel}: Vildt/kød skal vurderes af \textbf{uddannet person} (se §49.5)
        \item \textbf{Undtagelse}: Ikke krav ved direkte levering til endelig forbruger eller lokale detailbutikker (restauranter, vildtbutikker, supermarkeder)
        \item \textbf{Vurderingens indhold}:
          \begin{itemize}
            \item Adfærd før nedlæggelse (observeret af uddannet person eller beskrevet af nedlæggeren)
            \item Krop og organer efter nedlæggelse (tegn på sygdom, forurening m.m.)
          \end{itemize}
        \item \textbf{Jægererklæring} (kun storvildt til vildthåndteringsvirksomhed):
          \begin{itemize}
            \item Udstedes af uddannet person
            \item Indeholder: Dato, tid, sted for nedlæggelse; navn, underskrift, evt. jagttegnsnummer
            \item Kan dække flere dyr (med identifikationsnumre)
            \item To standarder: Normal (alt OK) eller unormal (problemer konstateret)
          \end{itemize}
        \item \textbf{Uden uddannet person til stede}: Hoved + organer (undt. gevir/horn/stødtænder + mave/tarme) skal medfølge; stadig sporbarhed (jæger, sted, dato)
        \item Kilder: Hygiejneforordningen bilag III afsnit IV kap. I pkt. 3; kap. II pkt. 4a
    \end{itemize}
\end{frame}